\subsection{Una agenda compartida de prevención familiar y ciudadana}
La combinación de fuentes propone un criterio de seguimiento: evaluar
avances por varias dimensiones, sin usar la caída de una tasa como evidencia
de protección universal. La disciplina observada en MICS, las actitudes de
género, la violencia letal y la victimización aportan señales diferentes.
Su articulación permite formular preguntas de política más precisas que una
clasificación global de seguridad o inseguridad.

La recomendación de Infosegura de dar seguimiento a avisos de violencia
intrafamiliar puede dialogar con intervenciones de crianza y prevención de
violencia de género. Se requiere información adicional para establecer
riesgos individuales y rutas de atención: la prevalencia nacional de MICS no
permite identificar personas que realizarán un aviso ni predecir un feminicidio.
Del mismo modo, una pregunta de justificación puede orientar diagnóstico de
normas, pero no debe emplearse para clasificar a una mujer como víctima o agresora.

\begin{table}[htbp]
\centering\small
\caption{Preguntas de política que surgen de la articulación}\label{tab:agenda}
\begin{tabularx}{\textwidth}{@{}YYY@{}}\toprule
Dimensión de MICS & Aporte de Infosegura & Pregunta y evidencia adicional requerida\\\midrule
Disciplina infantil & Enfoque de prevención y oportunidades de desarrollo & ¿Qué apoyos a cuidadores reducen prácticas violentas? Requiere evaluación y seguimiento propios.\\
Justificación de violencia & Seguimiento de avisos y muertes de mujeres & ¿Cómo se relacionan normas, acceso a servicios y respuesta institucional? Requiere datos vinculables y protección de identidad.\\
Unión temprana & Privaciones materiales y uso de espacios & ¿Qué restricciones de autonomía enfrentan las adolescentes? Requiere historia de unión, hogar de origen y trayectorias educativas.\\\bottomrule
\end{tabularx}
\fuente{Síntesis analítica del artículo y recomendaciones de Infosegura, diapositiva 13. Las preguntas no constituyen efectos estimados ni un sistema de predicción individual.}
\end{table}

Una futura investigación podría armonizar rondas de encuesta y territorios
para contrastar estas hipótesis. Tendría que documentar marcos muestrales,
cuestionarios, pesos y períodos. Con la evidencia presente, la recomendación
es conservar trazabilidad y complementar indicadores: no se identifica el
efecto del control del delito, de una reforma o de condiciones económicas
sobre las experiencias de niños y mujeres.
